\section{Introducción}
\label{sec:introduccion}

Cuando un grupo de personas trabaja de manera coordinada durante el tiempo suficiente,
algo notable ocurre: el grupo recuerda más de lo que cualquiera de sus miembros podría
recordar de forma individual.
Este fenómeno fue formalizado por Wegner \cite{wegner1987transactive} bajo el nombre de
\textbf{memoria transactiva}: un sistema cognitivo distribuido en el que cada individuo
mantiene no solo sus propios conocimientos, sino también un \emph{directorio} de quién
dentro del grupo es el experto en cada dominio.
Cuando un miembro necesita información fuera de su área, consulta ese directorio,
determina a quién preguntar y delega la recuperación.
La eficiencia del sistema surge precisamente de esa especialización emergente: el grupo
no duplica información, la distribuye, y la coordinación entre pares reemplaza
gradualmente a cualquier mediador central.

El paradigma de Wegner identifica tres mecanismos fundamentales.
Primero, la \emph{actualización del directorio}: tras cada interacción, todos los
miembros ---incluyendo quienes no fueron los expertos--- registran quién supo responder.
Segundo, la \emph{asignación de información}: la información nueva se canaliza
automáticamente hacia el miembro reconocido como especialista en ese dominio.
Tercero, la \emph{coordinación de recuperación}: las consultas futuras viajan directamente
al especialista sin pasar por un intermediario.
La transición de una \emph{fase temprana} mediada a una \emph{fase madura} punto a punto
es la marca distintiva de un sistema transactivo completamente formado.

Este trabajo plantea la siguiente pregunta de investigación:

\begin{quote}
¿Puede un sistema de memorias asociativas puras ---sin gradientes, sin
retropropagación y sin supervisión externa de las rutas--- desarrollar
dinámicamente el directorio transactivo de Wegner mediante la interacción
acumulativa con ejemplos de cada dominio?
\end{quote}

La respuesta que se propone y se verifica experimentalmente es afirmativa.
El sistema construido, denominado \textbf{MAE-TMS} (Multi-Agent Transactive Memory
System sobre Memorias Asociativas Entrópicas), se compone de tres agentes especialistas
entrenados sobre el conjunto de imágenes ETH-80 \cite{leibe2003eth80}: uno para la
categoría \textit{apple}, uno para \textit{horse} y uno para \textit{car}.
Un mediador central denominado TME (\textit{Transactive Memory Engine}) coordina la fase
temprana y registra en sus propios directorios qué agente responde a cada tipo de
consulta.
Conforme cada agente acumula registros en su memoria de directorio interna,
el TME se vuelve prescindible: en la fase madura, cualquier agente puede determinar
a quién reenviar una consulta desconocida sin intervención del mediador.

El sistema no emplea ningún mecanismo de optimización numérica.
Todas las operaciones ---registro, reconocimiento, recuperación y enrutamiento--- son
operaciones definidas formalmente sobre las \textbf{Memorias Asociativas Entrópicas}
(EAM) de Pineda \& Morales \cite{pineda2021entropic,pineda2024hetero}.
Las representaciones semánticas provienen de embeddings fastText \cite{bojanowski2017fasttext}
enriquecidos con relaciones de ConceptNet 5.7 \cite{speer2017conceptnet}, procesados con
spaCy \cite{honnibal2020spacy}; las representaciones visuales las produce un encoder
ResNet18 \cite{he2016resnet} preentrenado, cuya salida se cuantiza y almacena.

\subsection*{Contribuciones principales}

Este trabajo aporta tres contribuciones concretas:

\begin{enumerate}
    \item \textbf{Arquitectura de 14 memorias asociativas operando bidireccionalmente.}
    Cada agente emplea cuatro memorias (dos homo-asociativas de dominio y una
    hetero-asociativa de contenido texto$\leftrightarrow$imagen, más una hetero-asociativa
    de directorio); el TME incorpora dos directorios adicionales.
    Las memorias homo-asociativas actúan como árbitros de pertenencia y fuente de pesos
    para la proyección hetero-asociativa, resolviendo la asimetría fundamental entre
    dominios sin ninguna capa adicional.

    \item \textbf{Caracterización experimental completa del sistema.}
    Se reportan curvas de precisión temprana y madura en función del tamaño de la base
    de conocimiento $N$, la curva de capacidad (generalización sin vistas del cue), la
    formación del directorio (número de instancias necesarias para la especialización),
    el espacio semántico latente y la rejilla de recuperación imagen$\to$labels.

    \item \textbf{Hallazgo sobre el llenado teóricamente fiel.}
    Cuando las memorias se llenan siguiendo el esquema \emph{image-major} ---cada imagen
    registrada con cada label semántico de su clase, preservando la proporción natural
    de frecuencias--- los sesgos de masa quedan igualados por construcción (200 registros
    por agente) y el enrutamiento se torna directamente comparable entre agentes.
    Esto elimina la necesidad de calibraciones externas, con excepción de una única
    normalización por masa de registro (lectura B1) en el directorio, cuya justificación
    se deriva analíticamente de la definición de proyección hetero-asociativa.
\end{enumerate}

\subsection*{Estructura del documento}

El documento se organiza de la siguiente manera.
La Sección~\ref{sec:marco_teorico} presenta los fundamentos teóricos: la EAM de Pineda
\& Morales, la memoria hetero-asociativa 4D y el modelo de memoria transactiva de Wegner.
La Sección~\ref{sec:arquitectura} describe en detalle la arquitectura MAE-TMS, incluyendo
los diagramas de flujo del llenado, la fase temprana y la fase madura, así como la
formalización de todas las operaciones del sistema.
La Sección~4 detalla la metodología experimental.
La Sección~5 reporta los resultados.
La Sección~6 discute los hallazgos.
La Sección~7 concluye y señala trabajo futuro.
